% GENERATED FILE. Edit canonical lessons, then run npm run generate:books.
% canonical-source-hash: fnv1a64:09ed69786d206543
% canonical-lessons: ZH-C145-yuanyin, ZH-C145-wandian, ZH-C145-naozhong, ZH-C145-shuilanjiao, ZH-C145-gankuai

\chapter{Late Again}
\label{ch:zh-late-again}
\hlchaptermodality{145}

\begin{chapteropening}
\textbf{By the end of this chapter:} \emph{I can give reasons for being late: the train, the alarm clock, a lie-in.}

Complete the last lesson of chapter 145: I can give reasons for being late: the train, the alarm clock, a lie-in.
\end{chapteropening}

\section[yuányīn]{yuányīn — reason; cause}

\label{lesson:ZH-C145-yuanyin}



\begin{quote}
Before the new one: say the Mandarin for special, then the Mandarin for to find (that).
\end{quote}

\subsection*{You'll want to know: yuányīn}
\emph{yuányīn} — \textquotedblleft{}reason; cause\textquotedblright{}.

\begin{quote}
\emph{wǒ bù zhīdào shì shénme yuányīn} — I don't know what the reason is.
\end{quote}

Its characters come later; for now it is a word for the ear.

\begin{cousinweb}
\emph{yuányīn} is built from \emph{yuán} and \zh{因} 'because', both inherited from Middle Chinese, \emph{ngjwon} and \emph{'jin} (Guangyun, in Baxter's transcription). The dictionary glosses the first 'a broad plain' and the second 'to follow from'; how the first came to mean origin is not settled.
\end{cousinweb}

\begin{grammarlens}[title={What you've built}]
One more word for giving reasons.
\end{grammarlens}

\subsection*{Your turn}
\begin{itemize}
\raggedright
  \item \emph{Say it:} \emph{yuányīn}
  \item \emph{Say it:} \emph{yuányīn}, once more
  \item \emph{Say it:} \emph{wǒ bù zhīdào shì shénme yuányīn} — I don't know what the reason is.
  \item \emph{Recall:} say the Mandarin for special, then the Mandarin for to find (that), then say \emph{yuányīn} again
\end{itemize}

\begin{tcolorbox}[breakable,colback=teal!4,colframe=teal!35!black,title={Before you move on}]
What does yuányīn mean? (\textquotedblleft{}Reason; cause\textquotedblright{}.) Say it once more.
\end{tcolorbox}
\section[\texorpdfstring{\zh{晚点} (wǎndiǎn)}{晚点 (wǎndiǎn)}]{\zh{晚点} (wǎndiǎn) — to run behind schedule; to be late (of a train or plane)}

\label{lesson:ZH-C145-wandian}



\begin{quote}
Before the new one: say the Mandarin for to find (that), then the Mandarin for a reason.
\end{quote}

\subsection*{You'll want to know: \zh{晚点}}
\textbf{\zh{晚点}} — \emph{wǎndiǎn} — \textquotedblleft{}to run behind schedule; to be late (of a train or plane)\textquotedblright{}.

\begin{quote}
\textbf{\zh{火车晚点了半个小时}} \emph{huǒchē wǎndiǎn le bàn ge xiǎoshí} — The train was half an hour late.
\end{quote}

Two third tones: wǎndiǎn is said wándiǎn, the first rising.

\begin{cousinweb}
\textbf{\zh{晚点}} (\emph{wǎndiǎn}) is built from \zh{晚} and \zh{点}, both inherited from Middle Chinese, \emph{mjonX} and \emph{temX} (Guangyun, in Baxter's transcription), glossed 'dusk' and 'a dot'. Late against the hour is the sense of the pair, which the dictionary does not give.
\end{cousinweb}

\begin{grammarlens}[title={What you've built}]
One more word for giving reasons.
\end{grammarlens}

\subsection*{Your turn}
\begin{itemize}
\raggedright
  \item \emph{Say it:} \emph{wǎndiǎn}
  \item \emph{Say it:} \emph{wǎndiǎn}, once more
  \item \emph{Say it:} \emph{huǒchē wǎndiǎn le bàn ge xiǎoshí} — The train was half an hour late.
  \item \emph{Recall:} say the Mandarin for to find (that), then the Mandarin for a reason, then say \emph{wǎndiǎn} again
\end{itemize}

\begin{tcolorbox}[breakable,colback=teal!4,colframe=teal!35!black,title={Before you move on}]
What does \zh{晚点} mean? (\textquotedblleft{}To run behind schedule; to be late (of a train or plane)\textquotedblright{}.) Say it once more.
\end{tcolorbox}
\section[nàozhōng]{nàozhōng — alarm clock}

\label{lesson:ZH-C145-naozhong}



\begin{quote}
Before the new one: say the Mandarin for a reason, then the Mandarin for to run behind schedule.
\end{quote}

\subsection*{You'll want to know: nàozhōng}
\emph{nàozhōng} — \textquotedblleft{}alarm clock\textquotedblright{}.

\begin{quote}
\emph{wǒ méiyǒu nàozhōng suǒyǐ wǒ chídào le} — I haven't got an alarm clock, so I was late.
\end{quote}

Its characters come later; for now it is a word for the ear.

\begin{cousinweb}
\emph{nàozhōng} is built from \emph{nào} 'to stir up', the \emph{nao} of \emph{rènao}, and \emph{zhōng} 'bell', the \emph{zhōng} of \emph{fēnzhōng}, both inherited from Middle Chinese, \emph{nraewH} and \emph{tsyowng} (Guangyun, in Baxter's transcription): a bell that stirs you up.
\end{cousinweb}

\begin{grammarlens}[title={What you've built}]
One more word for giving reasons.
\end{grammarlens}

\subsection*{Your turn}
\begin{itemize}
\raggedright
  \item \emph{Say it:} \emph{nàozhōng}
  \item \emph{Say it:} \emph{nàozhōng}, once more
  \item \emph{Say it:} \emph{wǒ méiyǒu nàozhōng suǒyǐ wǒ chídào le} — I haven't got an alarm clock, so I was late.
  \item \emph{Recall:} say the Mandarin for a reason, then the Mandarin for to run behind schedule, then say \emph{nàozhōng} again
\end{itemize}

\begin{tcolorbox}[breakable,colback=teal!4,colframe=teal!35!black,title={Before you move on}]
What does nàozhōng mean? (\textquotedblleft{}Alarm clock\textquotedblright{}.) Say it once more.
\end{tcolorbox}
\section[shuì lǎnjiào]{shuì lǎnjiào — to sleep in; to lie in}

\label{lesson:ZH-C145-shuilanjiao}



\begin{quote}
Before the new one: say the Mandarin for to run behind schedule, then the Mandarin for alarm clock.
\end{quote}

\subsection*{You'll want to know: shuì lǎnjiào}
\emph{shuì lǎnjiào} — \textquotedblleft{}to sleep in; to lie in\textquotedblright{}.

\begin{quote}
\emph{jīntiān shì xīngqītiān wǒ xiǎng shuì lǎnjiào} — It's Sunday, so I want a lie-in.
\end{quote}

Its characters come later; for now it is a word for the ear.

\begin{cousinweb}
\emph{shuì lǎnjiào} is built from \emph{shuì} 'to sleep', \emph{lǎn} 'lazy' and \emph{jiào} 'a sleep', all inherited from Middle Chinese, \emph{dzyweH}, \emph{lanX} and \emph{kaewH} (Guangyun, in Baxter's transcription); the dictionary glosses the last with the word \emph{shuìjiào} itself.
\end{cousinweb}

\begin{grammarlens}[title={What you've built}]
One more word for giving reasons.
\end{grammarlens}

\subsection*{Your turn}
\begin{itemize}
\raggedright
  \item \emph{Say it:} \emph{shuì lǎnjiào}
  \item \emph{Say it:} \emph{shuì lǎnjiào}, once more
  \item \emph{Say it:} \emph{jīntiān shì xīngqītiān wǒ xiǎng shuì lǎnjiào} — It's Sunday, so I want a lie-in.
  \item \emph{Recall:} say the Mandarin for to run behind schedule, then the Mandarin for alarm clock, then say \emph{shuì lǎnjiào} again
\end{itemize}

\begin{tcolorbox}[breakable,colback=teal!4,colframe=teal!35!black,title={Before you move on}]
What does shuì lǎnjiào mean? (\textquotedblleft{}To sleep in; to lie in\textquotedblright{}.) Say it once more.
\end{tcolorbox}
\section[gǎnkuài]{gǎnkuài — quickly; hurry}

\label{lesson:ZH-C145-gankuai}



\begin{quote}
Before the new one: say the Mandarin for alarm clock, then the Mandarin for to lie in.
\end{quote}

\subsection*{You'll want to know: gǎnkuài}
\emph{gǎnkuài} — \textquotedblleft{}quickly; hurry\textquotedblright{}.

\begin{quote}
\emph{bā diǎn shàngkè wǒmen gǎnkuài zǒu} — Class is at eight, so let's hurry.
\end{quote}

Its characters come later; for now it is a word for the ear.

\begin{cousinweb}
\emph{gǎnkuài} is built from \emph{gǎn} 'to hurry' and \zh{快} 'fast'. \zh{快} is inherited from Middle Chinese \emph{khwaejH} (Guangyun, in Baxter's transcription); no dictionary we use records an older form of the first, so its history is not settled.
\end{cousinweb}

\begin{grammarlens}[title={What you've built}]
One more word for giving reasons.
\end{grammarlens}

\subsection*{Your turn}
\begin{itemize}
\raggedright
  \item \emph{Say it:} \emph{gǎnkuài}
  \item \emph{Say it:} \emph{gǎnkuài}, once more
  \item \emph{Say it:} \emph{bā diǎn shàngkè wǒmen gǎnkuài zǒu} — Class is at eight, so let's hurry.
  \item \emph{Recall:} say the Mandarin for alarm clock, then the Mandarin for to lie in, then say \emph{gǎnkuài} again
\end{itemize}

\begin{tcolorbox}[breakable,colback=teal!4,colframe=teal!35!black,title={Before you move on}]
What does gǎnkuài mean? (\textquotedblleft{}Quickly; hurry\textquotedblright{}.) Say it once more.
\end{tcolorbox}
